\begin{vidu}
Phương trình dao động điều hoà là $x=5\sin\!\left(2\pi t+\dfrac{\pi}{3}\right)\ (cm)$.  
Điền kết quả vào các ô trống trong bảng sau:	
\begin{center}
\renewcommand{\arraystretch}{1.3}
\begin{tabular}{|p{10cm}|p{3cm}|}
\hline Tần số của dao động là & \\ \hline
\hline Chu kì của dao động là & \\ \hline
\hline Pha ban đầu của dao động là & \\ \hline
\hline
Số dao động vật thực hiện trong 5 s là & \\ \hline
 \hline Chiều dài quỹ đạo của dao động là & \\ \hline
 \hline Quãng đường vật đi được sau 2 dao động là & \\ \hline
 \hline
 Pha của vật khi $t=6\ s$ là& \\ \hline
 \hline
 Li độ của vật khi $t=6\ s$ là& \\ \hline
\end{tabular}
\end{center}

\loigiai{
\begin{center}
\renewcommand{\arraystretch}{1.3}
\begin{tabular}{|p{10cm}|p{3cm}|}
\hline Tần số của dao động là & $1\ Hz$ \\ \hline
Chu kì của dao động là & $1\ s$ \\ \hline
Pha ban đầu của dao động là & $-\dfrac{\pi}{6}\ rad$ \\ \hline
Số dao động vật thực hiện trong $5\ s$ là & $5$ \\ \hline
Chiều dài quỹ đạo của dao động là & $10\ cm$ \\ \hline
Quãng đường vật đi được sau $2$ dao động là & $40\ cm$ \\ \hline
Pha của vật khi $t = 6\ s$ là & $\dfrac{71\pi}{6}\ rad$ \\ \hline
Li độ của vật khi $t = 6\ s$ là & $2,5\sqrt{3}\ cm$ \\ \hline
\end{tabular}
\end{center}
}
\end{vidu}
%\dotlineEX{6}
\begin{vidu}
Xác định biên độ dao động $A$, pha ban đầu và pha ở thời điểm t của các dao động có phương trình sau:
\begin{center} 
\begin{tabular}{| >{\rule[-10pt]{0pt}{30pt}} c | >{\rule[-10pt]{0pt}{30pt}} c | >{\rule[-10pt]{0pt}{30pt}} c | >{\rule[-10pt]{0pt}{30pt}} c |}		\hline
Đề bài & Biên độ (cm) & Pha ban đầu (rad) &\hspace{1cm} Pha (rad)\hspace{1cm} \\
\hline
$x=-4 \cos \left(10 \pi t-\dfrac{\pi}{4}\right)\ cm$ & & & \\
\hline
\hline
$x=4 \cos \left(-5 \pi t+\dfrac{\pi}{6}\right)\ cm$ & & & \\
\hline
\hline
$x=-6 \cos \left(-5 \pi t+\dfrac{\pi}{3}\right)\ cm$ & & & \\
\hline
\hline
$x=-\cos \left(2 \pi t-\dfrac{\pi}{12}\right)\ cm$ & & & \\
\hline
\hline
$x=4 \sin \left(10 \pi t-\dfrac{\pi}{4}\right)\ cm$ & & & \\
\hline
\hline
$x=5 \sin \left(-5 \pi t+\dfrac{\pi}{6}\right)\ cm$ & & & \\
\hline
\hline
$x=-3 \sin \left(\pi t-\dfrac{\pi}{8}\right)\ cm$ & & & \\
\hline
\hline
$x=-\sin \left(-4 \pi t-\dfrac{7 \pi}{12}\right)\ cm$ & & & \\
\hline
\end{tabular}
\end{center} 

\loigiai{
\begin{center}
\begin{tabular}{| >{\rule[-10pt]{0pt}{30pt}} c | >{\rule[-10pt]{0pt}{30pt}} c | >{\rule[-10pt]{0pt}{30pt}} c | >{\rule[-10pt]{0pt}{30pt}} c |}		\hline
Đề bài & Biên độ (cm) & Tần số góc (rad/s) &\hspace{1cm} Pha ban đầu (rad)\hspace{1cm} \\
\hline
$x=-4 \cos \left(10 \pi t-\dfrac{\pi}{4}\right)\ cm$ & $4$ & $10 \pi$ & $\dfrac{3\pi}{4}$ \\
\hline
$x=4 \cos \left(-5 \pi t+\dfrac{\pi}{6}\right)\ cm$ & $4$ & $5 \pi$ & $-\dfrac{\pi}{6}$ \\
\hline
$x=-6 \cos \left(-5 \pi t+\dfrac{\pi}{3}\right)\ cm$ & $6$ & $5 \pi$ & $\dfrac{4\pi}{3}$ \\
\hline
$x=-\cos \left(2 \pi t-\dfrac{\pi}{12}\right)\ cm$ & $1$ & $2 \pi$ & $\dfrac{11\pi}{12}$ \\
\hline
$x=4 \sin \left(10 \pi t-\dfrac{\pi}{4}\right)\ cm$ & $4$ & $10 \pi$ & $-\dfrac{3\pi}{4}$ \\
\hline
$x=5 \sin \left(-5 \pi t+\dfrac{\pi}{6}\right)\ cm$ & $5$ & $5 \pi$ & $-\dfrac{\pi}{3}$ \\
\hline
$x=-3 \sin \left(\pi t-\dfrac{\pi}{8}\right)\ cm$ & $3$ & $\pi$ & $\dfrac{7\pi}{8}$ \\
\hline
$x=-\sin \left(-4 \pi t-\dfrac{7 \pi}{12}\right)\ cm$ & $1$ & $4 \pi$ & $\dfrac{5\pi}{12}$ \\
\hline
\end{tabular}
\end{center}
}
\end{vidu}
%\dotlineEX{17}
\begin{vidu}
Một vật dao động có đồ thị li độ – thời gian được mô tả như hình vẽ.  
\begin{center} 
% Giữ nguyên hình vẽ gốc
\begin{tikzpicture}[draw=Mapcolor, very thick,scale=1,transform shape]
%Hệ trục toạ độ
\tkzDefPoints{0/0/x', 10/0/x, 0/-2.2/y', 0/2.2/y}
\tkzDrawSegments[Mapcolor, very thick, ->,add = 0 and .05](x',x)
\tkzDrawSegments[Mapcolor, very thick, ->,add = 0 and .05](y',y)
\tkzLabelPoint[below](x){$t\ (s)$}
\tkzLabelPoint[above=10pt](y){$x\ (cm)$}
%Chia khoảng trên trục x, y
\foreach  \x  in  {1,2,...,9}
{\draw[thin]  (\x,-0.1)  --  (\x,0.1);}
\foreach  \y  in  {-2,-1,1,2}
{\draw[thin]  (-0.1,\y)  --  (0.1,\y);}
%Đánh số toạ độ
\pgfkeys{/pgf/number format/use comma}
\foreach \x[evaluate=\x as \text using \x*0.1] in {1,2,...,9}{
\node[below=3pt] at (\x,0){\pgfmathprintnumber[fixed,fixed zerofill,precision=1]{\text}};}
\foreach \y[evaluate=\y as \text using \y*0.1] in {-2,-1,1,2}{
\node[left=3pt] at (0,\y){\pgfmathprintnumber[fixed,fixed zerofill,precision=1]{\text}};}
\node[below left] at (0,0){$O$};
%Đồ thị
\def\A{2} % Biên độ
\def\T{4} % Chu kì
\def\Phi{-90} % Pha ban đầu
{\draw [domain=0:9, samples=500]%
plot (\x, {(\A)*cos(360/\T*\x+\Phi)});}
\draw[fill=white](4,0)circle(2pt)node[above left]{$C$};
\draw[dashed](3,-2)--(0,-2);
\draw[fill=white](3,-2)circle(2pt)node[above]{$B$};
\draw[dashed](2.33,-1)--(0,-1);
\draw[fill=white](2.33,-1)circle(2pt)node[above right]{$A$};
\end{tikzpicture}
\end{center}
Hãy điền kết quả vào các ô trống trong bảng dưới đây:

\begin{center}
\renewcommand{\arraystretch}{1.3}
\begin{tabular}{|p{10cm}|p{3.2cm}|}
\hline Biên độ dao động $A$ & \\ \hline
\hline
Chu kì dao động $T$ & \\ \hline
\hline
Tần số dao động $f$ & \\ \hline
\hline
Tần số góc $\omega$ & \\ \hline
\hline
Li độ tại điểm $A$ & \\ \hline
\hline
Li độ tại điểm $B$ & \\ \hline
\hline
Li độ tại điểm $C$ & \\ \hline
\end{tabular}
\end{center}



\loigiai{
\begin{center}
\renewcommand{\arraystretch}{1.3}
\begin{tabular}{|p{10cm}|p{3.2cm}|}
\hline Biên độ dao động $A$ & $0,2\ cm$ \\ \hline
Chu kì dao động $T$ & $0,4\ s$ \\ \hline
Tần số dao động $f$ & $2,5\ Hz$ \\ \hline
Tần số góc $\omega$ & $5\pi\ rad/s$ \\ \hline
Li độ tại điểm $A$ & $-0,1\ cm$ \\ \hline
Li độ tại điểm $B$ & $-0,2\ cm$ \\ \hline
Li độ tại điểm $C$ & $0\ cm$ \\ \hline
\end{tabular}
\end{center}
}
\end{vidu}
\begin{vidu}
Một con ong mật đang bay tại chỗ trong không trung, đập cánh với tần số khoảng $300\ Hz$.  
\begin{vdc}
Số dao động mà cánh ong mật thực hiện trong $1 \mathrm{~s}$ là bao nhiêu?
\shortans[oly]{300}
\end{vdc}
%\dotlineEX{2}
\begin{vdc}
Chu kì dao động của cánh ong là bao nhiêu mili giây (làm tròn kết quả đến chữ số hàng phần trăm)?
\shortans[oly]{3,33}	
\end{vdc}
%\dotlineEX{2}
\loigiai{
\begin{itemize}
\item Tần số $f$ cho biết số dao động mà vật thực hiện được trong $1 \mathrm{~s}$.
 \item Mà $f=300 \mathrm{~Hz}$ nên số dao động mà cánh ong mật thực hiện trong $1 \mathrm{~s}$ là 300 dao động. 
 \item Chu kì $T=\dfrac{1}{f}=\dfrac{1}{300} \mathrm{~s}$.
\end{itemize}
}
\end{vidu} \haidong{20}


%\loai{Chiều dài quỹ đạo}
\begin{vidu}%[3P1-1.2Y][Loai1]
Một vật nhỏ dao động điều hòa với biên độ 3 cm. Vật dao động trên đoạn thẳng dài
\choice
{12 cm}
{9 cm}
{3 cm}
{\True 6 cm}
\loigiai{
Quỹ đạo chuyển động của vật là $L = 2A = 2.3 = 6$ cm.
}
\end{vidu} 
%\dotlineEX{2}
%\loai{Quãng đường đi trong một chu kì}
\begin{vidu}
Một chất điểm dao động điều hoà trong 10 dao động toàn phần đi được quãng đường dài 120 cm. Quỹ đạo của dao động có chiều dài là
\choice
{\True 6 cm}
{12 cm}
{3 cm}
{9 cm}
%\haidong{2}
\loigiai{
}
\end{vidu} %\dotlineEX{2}
%\loai{Pha của dao động}
\begin{vidu}%[3P1-1.2B][Loai1]
Một vật nhỏ dao động điều hòa theo phương trình $x = A\cos 10t$ (t tính bằng s), A là biên độ. Tại $t = 2$ s, pha của dao động là
\choice
{5 rad}
{40 rad}
{10 rad}
{\True 20 rad}
%\haidong{2}
\loigiai{
Tại $t = 2$ s, pha của dao động là: $\phi  = 10.2 = 20$ rad.
}
\end{vidu} %\dotlineEX{2}

%\loai{Xác định biên độ dao động và tần số góc}
\begin{vidu}%[3P1-1.2B][Loai1] 
Một vật dao động điều hoà theo phương trình $x = 3\cos(5\pi t + \dfrac{\pi}{6})$ cm. Biên độ dao động và tần số góc của vật là
\choice
{$-3\ cm$ và 5$\pi$ rad/s}
{3 cm và $-5\pi$ rad/s}
{\True 3 cm và 5$\pi$ rad/s}
{$3\pi$ cm và 5$\pi$ rad/s}
%\haidong{2}
\loigiai{
Ta có: $x=A\cos \left(\omega t+\varphi\right) =3\cos(5\pi t + \dfrac{\pi}{6})\ cm$
$\Rightarrow A = 3\ cm,\ \omega = 5\pi\ rad/s$}
\end{vidu} %\dotlineEX{3}

%\loai{Tính chu kì và tần số}
\begin{vidu} %[3P1-1.2B][Loai1] 
Một chất điểm dao động điều hòa với phương trình $x = 2\cos(\pi t + \dfrac{\pi}{3})$ cm. Chu kì và tần số của dao động là
\choice
{0,5 s; 2 Hz}
{\True 2 s; 0,5 Hz}
{2 s, $\pi$ Hz}
{0,5 s, $\pi$ rad/s}
%\haidong{2}
\loigiai{
\begin{itemize}
\item Chu kì $T = \dfrac{2\pi}{\pi} = 2\ s$.
\item Tần số $f = \dfrac{1}{T} = 0,5\ Hz$.
\end{itemize}}

\end{vidu} %\dotlineEX{3}

%\loai{Xác định pha ban đầu của hàm cos}
\begin{vidu}%[3P1-1.2B][Loai1] 
Một chất điểm dao động điều hòa với phương trình $x = 4\cos(10\pi t + \pi)$ cm. Pha ban đầu của chất điểm là
\choice
{\True $\pi$ rad}
{$0$ rad}
{$\dfrac{\pi}{2}$ rad}
{$10\pi t + \pi$ rad}
%\haidong{2}
\loigiai{
Pha ban đầu của chất điểm là $\pi$ rad.
}
\end{vidu} %\dotlineEX{2}

%\loai{Xác định pha ban đầu của hàm sin}
\begin{vidu} %[3P1-1.2B][Loai1]  
Một chất điểm dao động điều hòa với phương trình $x = 2\sin\left(\pi t + \dfrac{\pi}{2}\right)$ cm. Pha ban đầu của dao động trên là
\choice
{\True $0$ rad}
{$\dfrac{\pi}{2}$ rad}
{$3\dfrac{\pi}{2}$ rad}
{$\pi$ rad}
%\haidong{2}
\loigiai{
Ta có:
$x = 2\sin\left(\pi.t+ \dfrac{\pi}{2}\right)\ cm= 2cos\left(\pi.t+ \dfrac{\pi}{2}- \dfrac{\pi}{2}\right)\ cm = 2cos\left(\pi.t\right)$ cm.\\
Pha ban đầu của dao động trên là 0 rad.
}
\end{vidu} %\dotlineEX{2}



%\loai{Xác định pha ban đầu khi \bth{x=A\cos(\omega t+\varphi)} với $\varphi>\pi$ hoặc $\varphi<-\pi$}
\begin{vidu} %[3P1-1.2B][Loai1]  
Một dao động điều hòa với phương trình $x = 5\cos\left(2\pi t - \dfrac{3\pi}{2}\right)$ cm. Pha ban đầu của dao động của vật là
\choice
{$-\dfrac{\pi}{2}$ rad}
{\True $\dfrac{\pi}{2}$ rad}
{$\dfrac{3\pi}{2}$ rad}
{$-\pi$ rad}
%\haidong{2}
\loigiai{
Ta có:
$x = 5\cos\left(2\pi.t- \dfrac{3\pi}{2}\right)\ cm= 5\cos\left(2\pi.t- \dfrac{3\pi}{2}+2\pi\right)\ cm= 5\cos\left(2\pi.t+ \dfrac{\pi}{2}\right)$ cm.\\
Pha ban đầu của dao động trên là $\dfrac{\pi}{2}$ rad.
}
\end{vidu} %\dotlineEX{2}

